\documentclass[11pt,a4paper]{article}
\usepackage[margin=23mm]{geometry}
\usepackage{amsmath,amssymb,booktabs,tabularx,array,hyperref,fancyhdr,listings}
\usepackage[T1]{fontenc}
\usepackage{lmodern}
\hypersetup{colorlinks=true,urlcolor=blue,linkcolor=blue}
\pagestyle{fancy}\fancyhf{}\fancyhead[L]{\small Wenu-imaging}
\fancyhead[R]{\small Numerical interface review / 3 November 2026}
\fancyfoot[C]{\thepage}
\setlength{\headheight}{15pt}
\setlength{\parindent}{0pt}\setlength{\parskip}{6pt}
\setlength{\emergencystretch}{3em}
\lstset{basicstyle=\ttfamily\small,breaklines=true,columns=fullflexible}
\begin{document}
{\LARGE\bfseries Joint shutter and linearity fit}\\[5pt]
{\large Proposed numerical interfaces for review}

\textbf{Status.} The paired gain/RON helper at
\texttt{e0b16476} passed the actual Intel Mac C17 gate: two tests, 2.58 seconds,
all three preservation checks true. Evidence:
\path{/Users/fselman/Software/build/wenu-imaging/mac-gain-pair-AdGQvwqZ}.
Issue 228 (the corrected synthetic test expectation) is closed.
This note proposes the next numerical interfaces; no joint fitter is implemented.

The 2 November detector-characterization formula guide remains controlling.
This review supplies the numerical contract and explicitly identifies the missing
Eclipse functionality. Pixel nonlinearity thresholds, lamp-drift fitting and
real-data execution are not selected here.

\section*{1. Scientific model and symbols}
For a directly modeled signal in region $r$, exposure $j$ and amplifier $a$:
\begin{align}
Q_a(s)&=s+c_{2,a}s^2+c_{3,a}s^3,\\
Q_a(s_{rj})&=L_r(t_j+\delta),\qquad t_j+\delta>0 .
\end{align}
There is \textbf{one signed $\delta$ for the whole mosaic}, separate positive
illumination rates $L_r$, and amplifier-specific $c_{2,a},c_{3,a}$.
There is no free additive signal offset to absorb bias error.
The shutter offset is determined jointly from the accepted measurements of
\textbf{all detectors}, with every detector using this same $\delta$.
Each ROI retains its own illumination rate and each amplifier its response
coefficients. Excluded or saturated samples contribute no information; usable
samples contribute through the full covariance-weighted objective in Section 2.
We do not fit separate detector delays and then average them.

\begin{tabularx}{\linewidth}{@{}lXl@{}}\toprule
Symbol&Meaning&Unit\\\midrule
$s_{rj}$&Measured bias-subtracted signal; see aggregation rule below&ADU\\
$Q_a(s)$&Linearized signal for amplifier $a$&ADU\\
$c_{2,a},c_{3,a}$&Quadratic and cubic inverse-response coefficients&
ADU$^{-1}$, ADU$^{-2}$\\
$L_r$&Linearized illumination rate in region $r$&ADU/s\\
$t_j$&Input duration in the explicitly established time convention&s\\
$\delta$&Common duration offset; positive means longer exposure&s\\
$\widehat s_i$&Predicted measured signal for observation $i$&ADU\\
$e_i=s_i-\widehat s_i$&Signal-space residual&ADU\\
$C$&Covariance of accepted signal measurements&ADU$^2$\\
$\theta$&Ordered vector of fitted coefficients, rates and common offset&mixed\\\bottomrule
\end{tabularx}

\textbf{Pixel/ROI distinction.} A nonlinear transform of a mean does not equal
the mean of transformed pixels. The evaluator must predict individual pixel
signals before aggregating over an ROI's accepted population. An ROI-mean
approximation is permitted only with an explicit validation bound over its
signal distribution and fitted interval. A table of ROI means alone does not
establish that bound; an unvalidated approximation must be refused.

Known BPM pixels are excluded from reference populations according to configured
bits. Saturation/nonfinite selection is per sample and independent of permanent
defect classification. A saturated central ROI does not exclude unsaturated
corner ROIs in that exposure. Original exposure order and pair identities remain.

\clearpage
\section*{2. Objective, validity and derivatives}
The prediction is obtained on the validated monotonic branch:
\begin{align}
\widehat s_{rj}
 &=Q_a^{-1}\!\left[L_r(t_j+\delta)\right],\\
Q'_a(s)&=1+2c_{2,a}s+3c_{3,a}s^2>0,\\
\chi^2(\theta)&=e(\theta)^{\mathsf T}C^{-1}e(\theta).
\end{align}
$C$ is explicitly supplied and held fixed for one optimization. Diagonal weights
are permitted only when justified. Shared master-bias errors, shared reference
measurements and overlapping ROIs must be represented or removed by using an
independent nonduplicated population. Do not silently replace $C$ by its diagonal.

For independent observations with known standard uncertainties $\sigma_i$:
\[
\chi^2=\sum_i\left(\frac{e_i}{\sigma_i}\right)^2 .
\]
A parameter-dependent noise model needs a separately reviewed likelihood and
weight-update policy; it is not silently substituted for this fixed-$C$ objective.
Unknown covariance cannot be relabeled as a measured uncertainty.

For a directly modeled signal, let $u=\widehat s_{rj}$ and $D=Q'_a(u)$.
Implicit differentiation gives:
\[
\frac{\partial u}{\partial\delta}=\frac{L_r}{D},
\qquad
\frac{\partial u}{\partial L_r}=\frac{t_j+\delta}{D},
\qquad
\frac{\partial u}{\partial c_{2,a}}=-\frac{u^2}{D},
\qquad
\frac{\partial u}{\partial c_{3,a}}=-\frac{u^3}{D}.
\]
The residual Jacobian is $J_{ik}=\partial e_i/\partial\theta_k$, the negative of
the applicable prediction derivative. For pixel-before-ROI prediction, aggregate
the prediction derivatives with exactly the same weights as the predicted
signals. These equations do not authorize treating broad ROI means as single pixels.

Candidate evaluation must reject nonpositive effective durations, nonpositive
illumination rates, nonmonotonic response on the full declared signal interval,
and targets outside that branch. No extrapolation or silent coefficient clipping.
The minimum of $Q'_a$ on an interval can be checked at its endpoints and its
interior quadratic vertex, when present. The numerical inversion must be
bracketed within that same interval, with a reported convergence tolerance.

\textbf{Identifiability.} Constrained bias subtraction is required: an unknown
bias pedestal can mimic shutter delay. Arbitrary independent lamp factors would
absorb exposure-time response and are not fitted. The ten-level OCAM sequence is
assessed by its usable coverage/rank, not rejected solely for having fewer than
the representative annual 22 pairs.

\section*{3. Numerical scaling and uncertainty}
Use positive caller-declared parameter scales $d_k$ and dimensionless coordinates
$\theta_k=d_k x_k$. Scales are numerical units, not scientific priors.
They must be recorded so that rank/conditioning decisions are reproducible.

If $C=HH^{\mathsf T}$, with $H$ a nonsingular lower triangular covariance
factor, define whitened residuals and Jacobian:
\[
\widetilde e=H^{-1}e,\qquad
\widetilde J=H^{-1}J,\qquad
A=\widetilde J\,\mathrm{diag}(d_k).
\]
The input contract supplies $H$ or an explicitly diagonal equivalent.
The recipe must establish its provenance; the numerical service checks
dimensions, finite entries and positive diagonal. No diagonal jitter is added
to hide exact duplicate observations or singular covariance.

\clearpage
\section*{4. Eclipse reuse and the proposed additions}
The reviewed source is the current gain/RON candidate, descended from accepted
main. Existing public interfaces are preserved.

\begin{tabularx}{\linewidth}{@{}p{36mm}XX@{}}\toprule
Existing resource&What it supplies&What is missing\\\midrule
\texttt{optimization.c: minimize}
&Simplex reflection, expansion and contraction scaffolding
&Checked context callback, parameter bounds/scales, finite/validity states and
explicit parameter/objective convergence controls\\
\texttt{matrix.c: matrix\_leastsq}
&Matrix storage and normal-equation least squares
&Pivoted QR, numerical rank and scaled condition reporting\\
\texttt{fit\_curve.c}
&Polynomial and simple response fits
&Shared signed mosaic timing with constrained validity and covariance\\
\texttt{detector.c}
&Legacy time-as-signal-polynomial fit; cubic correction primitive
&The accepted joint signal-space objective and complete result contract\\\bottomrule
\end{tabularx}

\textbf{Addition A: checked optimizer in the existing optimization module.}
Proposed entry point:
\begin{lstlisting}[language=C]
int minimize_checked(
    const optimization_problem *problem,
    const optimization_options *options,
    optimization_result **out);
void optimization_result_del(optimization_result *result);
\end{lstlisting}
The borrowed problem contains parameter count, initial simplex, positive scales,
finite bounds, a context pointer and a cost callback. The callback distinguishes
a valid finite cost, an inadmissible model point and a hard evaluation failure.
Options explicitly set evaluation limits, parameter-spread tolerance and
objective-spread tolerance. An initial feasible simplex is required.

Extend the existing simplex scaffolding; do not introduce another optimizer
inside a recipe. Invalid trial points are rejected rather than clipped.
Convergence requires feasible vertices and both declared spread checks.
Return an explicit nonconvergence state with diagnostics when limits are reached.
A converged simplex is a candidate minimum, not evidence of a unique or
scientifically identified solution. Repeated starts/profile checks belong to
the detector wrapper. Keep the old \texttt{minimize} behavior/API intact.

\textbf{Addition B: checked least-squares/QR diagnostic in the matrix module.}
\begin{lstlisting}[language=C]
int matrix_leastsq_qr_checked(
    const matrix *design, const matrix *rhs,
    const double *column_scales, double rank_tolerance,
    matrix_qr_result **out);
void matrix_qr_result_del(matrix_qr_result *result);
\end{lstlisting}
Matrices follow existing row-major storage and dimensions. The design is $M$
observations by $P$ parameters, with $M\ge P$; the first increment accepts one
$M$-element right-hand side. Positive column scales and a relative rank tolerance
are explicit. Use column-pivoted Householder QR on the scaled design,
\[
A\Pi=UR ,
\]
where $\Pi$ is the column permutation, $U$ has orthonormal columns and $R$ is
upper triangular. Report pivot order, rank, residual norm and a clearly labeled
scaled condition estimate; do not call a diagonal ratio an exact condition number.
No normal-equation inversion, hidden ridge term or substituted zero solution.

\clearpage
\section*{5. Detector wrapper and fit acceptance}
After the two numerical interfaces are reviewed, the detector-level entry point is:
\begin{lstlisting}[language=C]
int detector_joint_response_fit_checked(
    const detector_joint_problem *problem,
    const detector_joint_options *options,
    detector_joint_result **out);
void detector_joint_result_del(detector_joint_result *result);
\end{lstlisting}

The borrowed problem binds amplifier/ROI/exposure identities, signal units,
accepted samples, prediction/aggregation evaluator, fixed covariance factor and
parameter layout. Its evaluator must implement Section 1's pixel-before-ROI
contract or supply the validated restricted approximation. Compact measurement
blocks and bounded image access are permitted; do not retain all detector cubes
simultaneously. The wrapper supplies the shared signed offset, cubic model,
analytic derivatives and validity checks to the reusable numerical routines.

Options explicitly set parameter scales/bounds, accepted signal intervals,
initial simplexes, convergence/rank/conditioning limits and profile settings.
No numerical scientific defaults or automatic polynomial-degree selection
are introduced by this proposal.

The result owns fitted parameters with units/identities, full parameter
covariance when available, residuals, accepted/excluded measurement references,
valid signal intervals, rank/condition diagnostics, evaluation counts,
bound/convergence states and shutter-profile diagnostics.
Inputs remain immutable; general allocation/invalid-input errors leave an
occupied or caller-owned output unchanged. Dedicated destructors release every
result allocation. Numerical/scientific unavailability is explicit and is
not published as a valid zero shutter delay.

For known absolute covariance and a locally identifiable interior solution:
\[
\mathrm{Cov}(\theta)\simeq
\left(J^{\mathsf T}C^{-1}J\right)^{-1}.
\]
Compute this local linear approximation through QR/triangular solves, not by
forming the normal equations. Transform back from the dimensionless parameter
scales and label it as a local approximation. Do not multiply by reduced
$\chi^2$ when $C$ already expresses absolute measurement variance.
Relative-weight fitting, if requested later, needs its own scale/dof convention.

An optimizer stop is insufficient. Require adequate signal coverage, full rank,
acceptable scaled conditioning, valid monotonic branches and consistent
feasible solutions from declared starts. Profile the shutter term by holding
$\delta$ fixed and refitting the remaining parameters:
\[
\chi^2_{\rm profile}(\delta)=
\min_{\theta\,:\,\theta_\delta=\delta}\chi^2(\theta).
\]
A flat/unbounded profile, boundary-truncated interval or incompatible minima
means the shutter is undetermined. Store the profile grid and the declared
likelihood convention; do not manufacture a symmetric error bar.

\textbf{BPM classification remains separate.} A correctable amplifier curve does
not flag every pixel. A later per-pixel residual test may set reason bit 16 only
with adequate coverage and reviewed thresholds. Insufficient information is
coverage-unavailable, not a permanent defect. Known BPM pixels can still undergo
individual defect tests without joining the reference population.

\section*{6. Proposed implementation and verification order}
\begin{enumerate}
\item Checked QR and ownership/rank controls in Eclipse; independent small
matrix solutions, scaled/permuted columns, near-degeneracy, overflow and
allocation failures.
\item Checked extension of the existing optimizer; independently known minima,
bounds/inadmissible callbacks, parameter scaling, convergence limits and
ownership/failure checks. Preserve legacy callers with focused comparisons.
\item Detector wrapper and analytic derivative comparisons; recover negative,
zero and positive common shutter offsets and injected amplifier curvature.
Exercise degeneracy, shared-bias covariance, overlapping ROI treatment,
saturation, disconnected coverage and lamp-drift violations.
\end{enumerate}
Announce the purpose, synthetic inputs, numerical calls and pass criteria before
execution. Mac C17 first; representative Chapman validation follows for changed
code before full actual-data testing there. Reuse qualified gain/RON evidence;
do not rerun unchanged tests. No installation, real-data processing, merge or
branch cleanup follows from publishing this proposal.

\textbf{Review requested.} Approve or amend the two reusable numerical interfaces
and their contracts before implementation. The detector wrapper remains
conditional on an explicit pixel/ROI evaluator and covariance preparation
contract; this note does not claim those inputs already exist in the current BPM
service. The earlier accepted scientific model is unchanged.
\clearpage
\appendix
\section*{Teaching appendix: simplex and QR}
\textbf{Two different jobs.} Simplex searches for a minimum of the nonlinear
objective. QR solves linear least-squares problems and examines the local
independence of parameter effects. QR is not a second nonlinear optimizer.

\subsection*{A. Simplex: searching without derivatives}
For $P$ fitted parameters, a simplex consists of $P+1$ trial parameter vectors:
a triangle in two dimensions, a tetrahedron in three. Here ``simplex'' means
the Nelder--Mead downhill search, not the simplex algorithm of linear programming.
At each vertex we evaluate the same covariance-weighted residual cost
$\chi^2$. A lower cost means a better match to the measurements.

The method orders the vertices by cost and forms the centroid of the best $P$.
It reflects the worst vertex through that centroid. A promising reflection can
be extended by expansion; an unsuccessful one leads to contraction. If that
fails, the simplex shrinks toward its best vertex. These operations adapt the
location and size of the search without requiring analytic derivatives.

Parameter scales make timing, illumination and curvature steps comparable in
dimensionless coordinates. Bounds and physical checks reject inadmissible trial
points. Small vertex separation and small cost spread are stopping conditions,
not proof of a unique minimum. Repeated starts and the shutter profile are
therefore required. A nearly flat direction can leave the delay undetermined
even when the search stops successfully.

\subsection*{B. QR: separating independent parameter effects}
Near a fitted point, the residual changes approximately linearly with a small
parameter change. After covariance whitening and parameter scaling, write
\[
\widetilde e(x+\Delta x)\simeq\widetilde e(x)+A\Delta x,
\qquad \min_{\Delta x}\|A\Delta x-b\|^2,
\quad b=-\widetilde e(x).
\]
Here $x$ denotes dimensionless parameters, $\Delta x$ their small change,
$A$ the scaled, whitened residual Jacobian and $b$ the required residual change.
The double bars denote Euclidean length. Householder reflections rotate this
linear problem without changing lengths. With column pivoting,
\[
A\Pi=UR,\qquad Rz=U^{\mathsf T}b,\qquad \Delta x=\Pi z.
\]
$\Pi$ reorders parameter columns; $U$ has orthonormal columns; $R$ is upper
triangular; $z$ is the change in the reordered coordinates. For full column
rank, back substitution solves the triangular system. Pivoting helps expose
nearly dependent parameter effects; the declared tolerance determines numerical
rank. QR avoids forming $A^{\mathsf T}A$, whose condition number is the square
of that of $A$ for full column rank in the Euclidean norm.

Here QR checks whether shutter delay, illumination and curvature can be
distinguished locally and supports covariance calculations. It cannot create
missing signal coverage or resolve a physical degeneracy. A rank-deficient
result must be reported as unavailable, not converted into a precise delay.
\end{document}

